% GENERATED FILE. Edit canonical lessons, then run npm run generate:books.
% canonical-source-hash: fnv1a64:97b9bcba8a66d534
% canonical-lessons: PA-C281-program, PA-C281-irada, PA-C281-salah, PA-C281-taalna, PA-C281-turant

\chapter{Plans for the Holidays}
\label{ch:pa-a2-plans-for-the-holidays}
\hlchaptermodality{281}

\begin{chapteropening}
\textbf{By the end of this chapter:} \emph{I can talk about a plan for the coming holidays.}

Complete the last lesson of chapter 281: I can talk about a plan for the coming holidays.
\end{chapteropening}

\section[\texorpdfstring{\pa{ਪ੍ਰੋਗਰਾਮ} (progrām)}{ਪ੍ਰੋਗਰਾਮ (progrām)}]{\pa{ਪ੍ਰੋਗਰਾਮ} (progrām) — a plan; a programme}

\label{lesson:PA-C281-program}



\begin{quote}
Before the new one: say the Punjabi for to take out, then the Punjabi for perhaps.
\end{quote}

\subsection*{You'll want to know: \pa{ਪ੍ਰੋਗਰਾਮ}}
\textbf{\pa{ਪ੍ਰੋਗਰਾਮ}} — \emph{progrām} — \textquotedblleft{}a plan; a programme\textquotedblright{}.

\begin{quote}
\pa{ਛੁੱਟੀ} \pa{ਦਾ} \pa{ਪ੍ਰੋਗਰਾਮ} \pa{ਕੀ} \pa{ਹੈ}?
\end{quote}

\begin{quote}
\emph{chhuṭṭī dā progrām kī hai?}
\end{quote}

\begin{quote}
What is the plan for the holiday?
\end{quote}

A masculine noun.

\begin{cousinweb}
\textbf{\pa{ਪ੍ਰੋਗਰਾਮ}} (\emph{progrām}) is borrowed from English \emph{programme}.
\end{cousinweb}

\begin{grammarlens}[title={What you've built}]
One more word for saying what comes next.
\end{grammarlens}

\subsection*{Your turn}
\begin{itemize}
\raggedright
  \item \emph{Say it:} \emph{progrām}
  \item \emph{Say it:} \emph{progrām}, once more
  \item \emph{Say it:} \emph{chhuṭṭī dā progrām kī hai?} — What is the plan for the holiday?
  \item \emph{Recall:} say the Punjabi for to take out, then the Punjabi for perhaps, then say \emph{progrām} again
\end{itemize}

\begin{tcolorbox}[breakable,colback=teal!4,colframe=teal!35!black,title={Before you move on}]
What does \pa{ਪ੍ਰੋਗਰਾਮ} mean? (\textquotedblleft{}A plan; a programme\textquotedblright{}.) Say it once more.
\end{tcolorbox}
\section[\texorpdfstring{\pa{ਇਰਾਦਾ} (irādā)}{ਇਰਾਦਾ (irādā)}]{\pa{ਇਰਾਦਾ} (irādā) — an intention}

\label{lesson:PA-C281-irada}



\begin{quote}
Before the new one: say the Punjabi for perhaps, then the Punjabi for a plan.
\end{quote}

\subsection*{You'll want to know: \pa{ਇਰਾਦਾ}}
\textbf{\pa{ਇਰਾਦਾ}} — \emph{irādā} — \textquotedblleft{}an intention\textquotedblright{}.

\begin{quote}
\pa{ਮੇਰਾ} \pa{ਇਰਾਦਾ} \pa{ਪੱਕਾ} \pa{ਹੈ।}
\end{quote}

\begin{quote}
\emph{merā irādā pakkā hai.}
\end{quote}

\begin{quote}
My intention is firm.
\end{quote}

A masculine noun.

\begin{cousinweb}
\textbf{\pa{ਇਰਾਦਾ}} (\emph{irādā}) is a loan, usually traced through Persian to Arabic \emph{irāda}, \textquotedblleft{}will, wish\textquotedblright{}.
\end{cousinweb}

\begin{grammarlens}[title={What you've built}]
One more word for saying what comes next.
\end{grammarlens}

\subsection*{Your turn}
\begin{itemize}
\raggedright
  \item \emph{Say it:} \emph{irādā}
  \item \emph{Say it:} \emph{irādā}, once more
  \item \emph{Say it:} \emph{merā irādā pakkā hai.} — My intention is firm.
  \item \emph{Recall:} say the Punjabi for perhaps, then the Punjabi for a plan, then say \emph{irādā} again
  \item \emph{Write it:} copy \textbf{\pa{ਇਰਾਦਾ}} once, checking the whole word
\end{itemize}

\begin{tcolorbox}[breakable,colback=teal!4,colframe=teal!35!black,title={Before you move on}]
What does \pa{ਇਰਾਦਾ} mean? (\textquotedblleft{}An intention\textquotedblright{}.) Say it once more.
\end{tcolorbox}
\section[\texorpdfstring{\pa{ਸਲਾਹ} (salāh)}{ਸਲਾਹ (salāh)}]{\pa{ਸਲਾਹ} (salāh) — advice; counsel}

\label{lesson:PA-C281-salah}



\begin{quote}
Before the new one: say the Punjabi for a plan, then the Punjabi for an intention.
\end{quote}

\subsection*{You'll want to know: \pa{ਸਲਾਹ}}
\textbf{\pa{ਸਲਾਹ}} — \emph{salāh} — \textquotedblleft{}advice; counsel\textquotedblright{}.

\begin{quote}
\pa{ਪਿਤਾ} \pa{ਜੀ} \pa{ਚੰਗੀ} \pa{ਸਲਾਹ} \pa{ਦਿੰਦੇ} \pa{ਹਨ।}
\end{quote}

\begin{quote}
\emph{pitā jī chaṅgī salāh dinde han.}
\end{quote}

\begin{quote}
Father gives good advice.
\end{quote}

A feminine noun.

\begin{cousinweb}
\textbf{\pa{ਸਲਾਹ}} (\emph{salāh}) is a loan, usually traced to Persian \emph{salāh}, \textquotedblleft{}counsel\textquotedblright{}, itself from Arabic.
\end{cousinweb}

\begin{grammarlens}[title={What you've built}]
One more word for saying what comes next.
\end{grammarlens}

\subsection*{Your turn}
\begin{itemize}
\raggedright
  \item \emph{Say it:} \emph{salāh}
  \item \emph{Say it:} \emph{salāh}, once more
  \item \emph{Say it:} \emph{pitā jī chaṅgī salāh dinde han.} — Father gives good advice.
  \item \emph{Recall:} say the Punjabi for a plan, then the Punjabi for an intention, then say \emph{salāh} again
  \item \emph{Write it:} copy \textbf{\pa{ਸਲਾਹ}} once, checking the whole word
\end{itemize}

\begin{tcolorbox}[breakable,colback=teal!4,colframe=teal!35!black,title={Before you move on}]
What does \pa{ਸਲਾਹ} mean? (\textquotedblleft{}Advice; counsel\textquotedblright{}.) Say it once more.
\end{tcolorbox}
\section[\texorpdfstring{\pa{ਟਾਲਣਾ} (ṭālṇā)}{ਟਾਲਣਾ (ṭālṇā)}]{\pa{ਟਾਲਣਾ} (ṭālṇā) — to put off; to postpone}

\label{lesson:PA-C281-taalna}



\begin{quote}
Before the new one: say the Punjabi for an intention, then the Punjabi for advice.
\end{quote}

\subsection*{You'll want to know: \pa{ਟਾਲਣਾ}}
\textbf{\pa{ਟਾਲਣਾ}} — \emph{ṭālṇā} — \textquotedblleft{}to put off; to postpone\textquotedblright{}.

\begin{quote}
\pa{ਅਮਨ} \pa{ਹਰ} \pa{ਕੰਮ} \pa{ਕੱਲ੍ਹ} \pa{ਉੱਤੇ} \pa{ਟਾਲਦਾ} \pa{ਹੈ।}
\end{quote}

\begin{quote}
\emph{Aman har kamm kall utte ṭāldā hai.}
\end{quote}

\begin{quote}
Aman puts every job off until tomorrow.
\end{quote}

Past (I …): \pa{ਮੈਂ} \pa{ਯਾਤਰਾ} \pa{ਟਾਲੀ}, \emph{maiṁ yātrā ṭālī}.

\begin{cousinweb}
\textbf{\pa{ਟਾਲਣਾ}} (\emph{ṭālṇā}) is inherited from Sanskrit \emph{ṭalati}, \textquotedblleft{}moves aside\textquotedblright{}, through its causative \emph{ṭālayati}; CDIAL lists Punjabi \emph{ṭālṇā}, \textquotedblleft{}to evade\textquotedblright{}, under it (CDIAL 5450).
\end{cousinweb}

\begin{grammarlens}[title={What you've built}]
One more word for saying what comes next.
\end{grammarlens}

\subsection*{Your turn}
\begin{itemize}
\raggedright
  \item \emph{Say it:} \emph{ṭālṇā}
  \item \emph{Say it:} \emph{ṭālṇā}, once more
  \item \emph{Say it:} \emph{Aman har kamm kall utte ṭāldā hai.} — Aman puts every job off until tomorrow.
  \item \emph{Recall:} say the Punjabi for an intention, then the Punjabi for advice, then say \emph{ṭālṇā} again
  \item \emph{Write it:} copy \textbf{\pa{ਟਾਲਣਾ}} once, checking the whole word
\end{itemize}

\begin{tcolorbox}[breakable,colback=teal!4,colframe=teal!35!black,title={Before you move on}]
What does \pa{ਟਾਲਣਾ} mean? (\textquotedblleft{}To put off; to postpone\textquotedblright{}.) Say it once more.
\end{tcolorbox}
\section[\texorpdfstring{\pa{ਤੁਰੰਤ} (turant)}{ਤੁਰੰਤ (turant)}]{\pa{ਤੁਰੰਤ} (turant) — at once; immediately}

\label{lesson:PA-C281-turant}



\begin{quote}
Before the new one: say the Punjabi for advice, then the Punjabi for to put off.
\end{quote}

\subsection*{You'll want to know: \pa{ਤੁਰੰਤ}}
\textbf{\pa{ਤੁਰੰਤ}} — \emph{turant} — \textquotedblleft{}at once; immediately\textquotedblright{}.

\begin{quote}
\pa{ਮੈਂ} \pa{ਤੁਰੰਤ} \pa{ਘਰ} \pa{ਜਾਵਾਂਗਾ।}
\end{quote}

\begin{quote}
\emph{maiṁ turant ghar jāvāṅgā.}
\end{quote}

\begin{quote}
I will go home at once.
\end{quote}

\begin{cousinweb}
\textbf{\pa{ਤੁਰੰਤ}} (\emph{turant}) goes back to Sanskrit \emph{turant}, \textquotedblleft{}hastening\textquotedblright{}, as a word taken from Hindi: CDIAL lists Punjabi \emph{turant}, \textquotedblleft{}quickly\textquotedblright{}, under it and marks it as borrowed from Hindi (CDIAL 5879).
\end{cousinweb}

\begin{grammarlens}[title={What you've built}]
One more word for saying what comes next.
\end{grammarlens}

\subsection*{Your turn}
\begin{itemize}
\raggedright
  \item \emph{Say it:} \emph{turant}
  \item \emph{Say it:} \emph{turant}, once more
  \item \emph{Say it:} \emph{maiṁ turant ghar jāvāṅgā.} — I will go home at once.
  \item \emph{Recall:} say the Punjabi for advice, then the Punjabi for to put off, then say \emph{turant} again
  \item \emph{Write it:} copy \textbf{\pa{ਤੁਰੰਤ}} once, checking the whole word
\end{itemize}

\begin{tcolorbox}[breakable,colback=teal!4,colframe=teal!35!black,title={Before you move on}]
What does \pa{ਤੁਰੰਤ} mean? (\textquotedblleft{}At once; immediately\textquotedblright{}.) Say it once more.
\end{tcolorbox}
